% ------------------------------------------------------------
\begin{frame}[fragile]{Common Query Anti-Patterns}
  \begin{block}{Patterns that prevent index use or waste resources}
\begin{lstlisting}[style=sqlstyle]
-- Anti-pattern 1: function on indexed column
WHERE UPPER(name) = 'MILLER'      -- index on name not used
-- Fix: store data in normalized case, or use a functional index

-- Anti-pattern 2: leading wildcard
WHERE name LIKE '%miller%'         -- full table scan
-- Fix: use FULLTEXT index for free-text search

-- Anti-pattern 3: implicit type conversion
WHERE id = '42'    -- id is INT, '42' is VARCHAR -> may skip index
-- Fix: match the data type: WHERE id = 42

-- Anti-pattern 4: SELECT * on large tables
SELECT * FROM orders;              -- pulls every column, every row
-- Fix: always select only the columns you actually need

-- Anti-pattern 5: correlated subquery in SELECT
SELECT name, (SELECT COUNT(*) FROM orders WHERE customer_id = c.id)
FROM customer c;                   -- executes subquery per row!
-- Fix: use LEFT JOIN + GROUP BY instead
\end{lstlisting}
  \end{block}
\end{frame}

% ------------------------------------------------------------
\begin{frame}[fragile]{The N+1 Query Problem}
  The \textbf{N+1 problem} is one of the most common performance pitfalls
  in applications that use ORMs or generate SQL in loops.

  \begin{block}{What it is}
\begin{lstlisting}[style=sqlstyle]
-- 1 query to get all orders
SELECT * FROM orders;   -- returns N rows

-- Then N queries, one per order, to get the customer:
SELECT * FROM customer WHERE id = 1;
SELECT * FROM customer WHERE id = 2;
-- ... N more queries
\end{lstlisting}
  \end{block}

  \begin{block}{The fix: one JOIN instead of N+1 queries}
\begin{lstlisting}[style=sqlstyle]
-- Single query returns everything needed
SELECT o.id, o.total, c.name
FROM   orders o
JOIN   customer c ON o.customer_id = c.id;
\end{lstlisting}
  \end{block}

  \begin{alertblock}{Why it matters}
    With 10,000 orders, N+1 generates 10,001 round-trips to the database.
    Even at 1\,ms per query that is 10 seconds. A single JOIN takes
    milliseconds.
  \end{alertblock}
\end{frame}

% ------------------------------------------------------------
\begin{frame}[fragile]{Query Optimization -- Practical Checklist}
  \begin{block}{Step-by-step approach}
    \begin{enumerate}
      \item \textbf{Run \texttt{EXPLAIN ANALYZE}} -- identify the slowest
            operation in the plan (high \texttt{rows}, \texttt{ALL} access,
            \texttt{Using filesort}).
      \item \textbf{Check for missing indexes} on \texttt{WHERE},
            \texttt{JOIN ON}, and \texttt{ORDER BY} columns.
      \item \textbf{Avoid functions on indexed columns} in \texttt{WHERE}
            clauses -- rewrite the predicate instead.
      \item \textbf{Select only needed columns} -- drop \texttt{SELECT *};
            use covering indexes where possible.
      \item \textbf{Push filters early} -- filter rows as close to the
            source table as possible, before joins.
      \item \textbf{Replace correlated subqueries} with \texttt{JOIN} or
            window functions.
      \item \textbf{Batch writes}: wrap many small \texttt{INSERT} statements
            in a single multi-row \texttt{INSERT} or a transaction.
      \item \textbf{Update statistics}: run \texttt{ANALYZE TABLE}
            periodically so the planner has fresh data.
    \end{enumerate}
  \end{block}
\end{frame}
